% §7 Discussion

\subsection{Synthesis}
\label{sec:disc:synthesis}

We asked whether unsupervised pre-inference signals can support appropriate multi-LLM routing before generation.
On ARC-Challenge, the answer is \textbf{partially}.

The scientific contribution is dimensional: routing need decomposes into \textbf{latent dimensions}---task difficulty, model uncertainty, model disagreement, escalation potential---each operationalized by a prefill statistic.
Before generation, multiple dimensions predict routing opportunity; they encode \textbf{different aspects of routing need}.
A calibrated routing policy does not fully exploit what is available.

\subsection{Why the dimensional view matters}
\label{sec:disc:failures}

Treating \textbf{model uncertainty} as a proxy for \textbf{escalation potential} fails: opportunity and too-hard queries are indistinguishable on $H_w$ yet separate on $\Delta m_{\mathrm{gain}}$.
Routing policies must respect the dimensional structure documented in \S\ref{sec:char:finding2}.

\subsection{Future dimensions}
\label{sec:disc:future}

This study operationalizes four dimensions from query text and single-step prefill logits.
Richer measurements---generation trajectories, hidden states, paraphrase stability---may support additional latent dimensions (appendix Table~\ref{tab:dim-taxonomy}) without changing the dimensional abstraction.

\subsection{Limitations}
\label{sec:disc:limits}

Binary Llama~3.2 pool; ARC-Challenge only; one operationalization per main dimension; demonstration policy calibrated on CALIB.
